% Source: https://artofproblemsolving.com/community/c3h2056207_counting_problem
% Stable identity: aops-topic-2056207-q01
A finishing order is an ordered sequence of nonempty groups of runners (first-place group, then second, and so on). Let $a_n$ be the number of finishing orders for $n$ runners, with $a_0=1$. Choosing the set of $k\ge1$ runners who tie for first and then finishing the remaining $n-k$ runners gives
\[a_n=\sum_{k=1}^{n}\dbinom nk a_{n-k}.\]
Hence
\begin{align*}
a_1&=1,\\
a_2&=2a_1+a_0=3,\\
a_3&=3a_2+3a_1+a_0=9+3+1=13,\\
a_4&=4a_3+6a_2+4a_1+a_0=52+18+4+1=75,\\
a_5&=5a_4+10a_3+10a_2+5a_1+a_0=375+130+30+5+1=541,\\
a_6&=6a_5+15a_4+20a_3+15a_2+6a_1+a_0=3246+1125+260+45+6+1=4683.
\end{align*}
(Equivalently, $a_6=\sum_k k!\,S(6,k)=1+62+540+1560+1800+720$, where $S(6,k)$ counts the ways to split the runners into $k$ unordered groups.)

There are $\boxed{4683}$ possible finishing orders.
